\section{Canonical Business Twin Data, Function Persistence, and Boardroom Context Architecture}
\label{sec:business-twin-function-persistence}

\subsection{Purpose and Non-Negotiable Principle}

This section defines the reusable persistence and context architecture that every
business function must follow. It applies to Finance, Marketing, Sales,
Operations, Support, Builder, Pilot, and any future function introduced into
AION. Its purpose is to prevent each function from creating its own disconnected
storage keys, document stores, context packets, or model memory assumptions.

The governing principle is:

\begin{quote}
The scoped Business Container is the persistent source of truth. The Business
Map is the structured living representation of known business facts. Function
ledgers record departmental discovery, plans, work, evidence, and results. The
File Cabinet is an index of pointers to Business Container artifacts. The
Boardroom receives a bounded, immutable projection of this truth. External AI
providers are not relied upon to remember the business between meetings.
\end{quote}

Local renderer variables, browser \texttt{localStorage}, session storage, and
provider conversation history may be used as caches or resumable drafts, but
must never become the only persistent copy of material business information.

\subsection{System of Record and Projection Boundaries}

The platform separates persistent truth from views and projections:

\begin{verbatim}
Business Container (persistent truth)
    |
    +-- Business Identity
    +-- Business Structure
    +-- Business Map
    +-- Function Models and Ledgers
    +-- Evidence and Artifact Records
    +-- Goal Engine State, Memory, Evidence, Loops and Outcomes
    |
    +--> Business Twin / Boardroom Context Envelope (read projection)
    |        |
    |        +--> identical packet to OpenAI
    |        +--> identical packet to Claude
    |        +--> identical packet to Gemini
    |        +--> identical packet to other Boardroom providers
    |
    +--> File Cabinet (index pointers only)
\end{verbatim}

The Boardroom must not become a second mutable truth store. It reads a projection,
deliberates, creates decisions and packages, and writes approved results back
through controlled Business Container mutation paths. Likewise, the File Cabinet
does not own the bytes or canonical JSON of a business artifact; it indexes the
artifact's canonical location.

\subsection{Canonical Runtime Locations}

The AION Business runtime root is currently defined by
\texttt{AIONBusinessPaths.ROOT} as:

\begin{verbatim}
.runtime/AION_BUSINESS/
\end{verbatim}

The important persistent locations are:

\begin{verbatim}
.runtime/AION_BUSINESS/
    business_containers/<business_id>/
    workflow_file_cabinets/<business_id>/tree.json
    container_bindings/<business_id>/
    workspaces/<business_id>.json
    topologies/<business_id>.json
    audit/<business_id>/
    tasks/<business_id>/
    learning/<business_id>/
\end{verbatim}

The path owner is:

\begin{verbatim}
backend/modules/aion_business/runtime/paths.py
\end{verbatim}

Canonical structured-container persistence is owned by:

\begin{verbatim}
backend/modules/aion_business/runtime/business_container_repository.py
backend/modules/aion_business/runtime/business_container_service.py
backend/modules/aion_business/contracts/business_containers.py
\end{verbatim}

The Business Container repository writes one canonical JSON file per registered
container kind under the business-scoped directory.

\subsection{Canonical Business Identity}

Every record, file, model, ledger, receipt, Boardroom packet, and File Cabinet
pointer must resolve to one canonical \texttt{business\_id}. No function may
independently slug or rename the business identifier.

At the time of inspection, both \texttt{home\_fixed} and \texttt{home-fixed}
normalisation rules existed in different paths. This can create two separate
Business Containers and two separate File Cabinets. The platform must therefore
provide one canonical resolver used by frontend and backend code.

The canonical resolver must:

\begin{enumerate}
  \item accept the registered business identifier rather than the display name;
  \item reject global, root, temporary, and session-only targets;
  \item apply one documented slug rule exactly once;
  \item return the same identifier for Business Container, File Cabinet,
  Boardroom, artifact, audit, task, and function-ledger operations;
  \item support a controlled alias migration for legacy identifiers; and
  \item never silently fall back to a demo business.
\end{enumerate}

A representative identity object is:

\begin{verbatim}
{
  "business_id": "home-fixed",
  "display_name": "Home Fixed",
  "workspace_id": "home-fixed",
  "container_id": "home-fixed",
  "legacy_aliases": ["home_fixed"],
  "identity_status": "registered"
}
\end{verbatim}

All write endpoints must reject \texttt{business\_not\_registered},
\texttt{global}, \texttt{root}, \texttt{tmp}, \texttt{temp}, and
\texttt{session\_vfs} as persistent business targets.

\subsection{Existing Canonical Container Contracts}

The existing canonical container kinds include:

\begin{itemize}
  \item \texttt{business\_identity};
  \item \texttt{business\_structure};
  \item \texttt{brand\_foundation};
  \item \texttt{boardroom\_snapshot};
  \item \texttt{operational\_runtime\_summary};
  \item \texttt{goal\_engine\_state};
  \item \texttt{goal\_engine\_memory};
  \item \texttt{goal\_engine\_evidence};
  \item \texttt{goal\_engine\_experiments};
  \item \texttt{goal\_engine\_loops}; and
  \item \texttt{goal\_engine\_outcomes}.
\end{itemize}

The Business Identity container holds legal and trading names, business type,
sector, stage, owner, location, timezone, currency, domain, website, contact
details, description, and source references.

The Business Structure container already provides lists for services, revenue
streams, cost items, payment terms, fulfilment processes, functions, teams,
human agents, and AI agents. It is useful for cross-function structure, but it is
not a replacement for a full Finance model or another specialist function model.

\subsection{Required Extension for Function Models}

Every substantive function requires a first-class canonical model, or an agreed
generic function-model container contract. Finance must not be forced into a
marketing summary field or stored only as an arbitrary artifact.

Two acceptable implementation approaches are:

\begin{enumerate}
  \item register first-class kinds such as
  \texttt{business\_financial\_model},
  \texttt{business\_sales\_model},
  \texttt{business\_operations\_model}, and
  \texttt{business\_support\_model}; or
  \item register one \texttt{business\_function\_models} container containing
  schema-versioned records keyed by function.
\end{enumerate}

The first-class approach provides clearer validation and is preferred when the
function has a stable specialist schema. Every model must include:

\begin{itemize}
  \item schema version;
  \item business and function identifiers;
  \item model status and lifecycle state;
  \item structured facts and calculated values;
  \item assumptions and unknowns;
  \item evidence references;
  \item provenance;
  \item validation state;
  \item created, updated, and effective timestamps;
  \item revision or supersession information; and
  \item completion and approval receipts.
\end{itemize}

\subsection{The Business Map as the Living Business Record}

The Foundation runtime currently creates a Business Map with the schema
\texttt{aion.business\_map.voice\_foundation.v0.1}. It contains facts,
assumptions, unanswered fields, and department discovery gaps. Website scanning
adds evidence-aware facts containing confidence, source URL, receipt, provider,
model, and timestamps.

This useful frontend packet must be promoted into a canonical backend contract.
The canonical Business Map should not be a flat overwriteable dictionary. It
should retain provenance and permit multiple observations about the same field.

A reusable fact contract is:

\begin{verbatim}
{
  "fact_id": "fact_finance_monthly_revenue_2026_06",
  "business_id": "home-fixed",
  "function": "finance",
  "field": "average_monthly_revenue",
  "category": "revenue",
  "value": 4200,
  "unit": "EUR",
  "period": "2026-06",
  "classification": "founder_provided_fact",
  "confidence": 0.88,
  "verification_status": "unverified",
  "source_ref": "finance_discovery_turn_04",
  "evidence_refs": [],
  "supersedes": null,
  "effective_at": "2026-06-30T23:59:59Z",
  "observed_at": "2026-07-13T00:00:00Z",
  "created_at": "2026-07-13T00:00:00Z"
}
\end{verbatim}

Every fact must be classified as one of:

\begin{itemize}
  \item founder-provided fact;
  \item connected-system fact;
  \item uploaded-document fact;
  \item website-derived fact;
  \item calculated value;
  \item Boardroom-approved decision;
  \item assumption;
  \item disputed value; or
  \item unknown.
\end{itemize}

Corrections should append a new observation or revision and mark the older record
as superseded. Silent destructive replacement should be avoided for material
business facts.

\subsection{Function Discovery State}

Each function may keep a resumable discovery draft, but it must follow a common
envelope:

\begin{verbatim}
{
  "schema_version": "aion.function_discovery.v1",
  "business_id": "home-fixed",
  "function": "finance",
  "status": "in_progress",
  "current_section": "revenue",
  "current_question_id": "revenue_monthly",
  "answers": {},
  "transcript": [],
  "facts": [],
  "assumptions": [],
  "unknowns": [],
  "evidence_refs": [],
  "integration_refs": [],
  "progress": {
    "completed_fields": 0,
    "total_fields": 0,
    "percentage": 0
  },
  "safety": {
    "external_actions_allowed": false,
    "external_writes_allowed": false,
    "payments_allowed": false,
    "approval_required": true
  },
  "created_at": "...",
  "updated_at": "..."
}
\end{verbatim}

The renderer may cache this packet locally for fast resume. Material changes must
also be committed through the backend into the active Business Container. The
backend copy is authoritative after successful acknowledgement.

\subsection{Department Intelligence Ledger}

The existing frontend ledger uses \texttt{aion.departmentIntelligence.v1}. Its
normalised entries already support:

\begin{verbatim}
{
  "department": "finance",
  "status": "needs_discovery",
  "confidence": "empty",
  "discovery": {},
  "plan": {},
  "tasks": [],
  "runs": [],
  "evidence": [],
  "receipts": [],
  "alerts": [],
  "results": {},
  "boardroom_summary": "",
  "updated_at": "..."
}
\end{verbatim}

This is the correct conceptual ledger but currently relies on browser storage.
It must become a persistent Business Container record. Browser storage should
remain a cache, with a recorded sync state such as \texttt{pending},
\texttt{synced}, \texttt{conflicted}, or \texttt{failed}.

The function model and Department Intelligence Ledger serve different purposes:

\begin{itemize}
  \item the function model contains detailed specialist truth;
  \item the ledger contains operational state, discovery completion, plans,
  tasks, evidence, receipts, results, alerts, and a bounded Boardroom summary;
  \item the Boardroom reads both through a projection; and
  \item the ledger should reference model revisions rather than duplicate the
  entire model.
\end{itemize}

\subsection{Artifact Persistence Contract}

The existing artifact contract is implemented in:

\begin{verbatim}
backend/services/aion_mission_mode/business_container_artifacts.py
backend/services/aion_mission_mode/department_artifact_persistence.py
\end{verbatim}

It correctly requires a business-scoped target, department sub-container,
artifact type, provenance, artifact hash, source-payload hash, receipt hash,
storage path, and replay locator. It rejects global and temporary targets.

Allowed department sub-containers include Marketing, Pilot, Sales, Finance,
Operations, Support, and Builder. Supported user-usable types include PDF,
spreadsheet, document, text, Markdown, JSON, report, evidence bundle, export,
receipt, and Boardroom session.

A canonical artifact record should contain:

\begin{verbatim}
{
  "artifact_id": "artifact_...",
  "target": {
    "business_id": "home-fixed",
    "business_container_id": "home-fixed",
    "sub_container": "finance",
    "artifact_type": "spreadsheet",
    "artifact_name": "cashflow-2026.xlsx"
  },
  "provenance": {
    "mission_id": "finance-onboarding",
    "mission_run_id": "run_...",
    "step_id": "finance-file-import",
    "tool_id": "finance-file-ingestion",
    "created_by": "user_upload"
  },
  "artifact_hash": "sha256:...",
  "source_payload_hash": "sha256:...",
  "receipt_hash": "sha256:...",
  "storage_path": "business_containers/home-fixed/finance/spreadsheets/...",
  "storage_state": "committed_to_business_container",
  "session_vfs_only": false,
  "live_side_effects_enabled": false,
  "replay_locator": {}
}
\end{verbatim}

At the time of inspection, the artifact commit helper builds and validates this
record but does not itself write the physical artifact. A production backend
commit endpoint must therefore:

\begin{enumerate}
  \item validate the canonical business identifier and sub-container;
  \item stream the uploaded bytes into a temporary business-scoped staging area;
  \item enforce size and file-type policy;
  \item calculate a content hash;
  \item reject path traversal and unsafe filenames;
  \item place the artifact under the canonical business/function path;
  \item write the artifact record and receipt atomically;
  \item return the canonical storage path, hash, receipt, and replay locator; and
  \item only then create or update the File Cabinet pointer.
\end{enumerate}

No UI should display \texttt{committed\_to\_business\_container} before the
backend has confirmed that the bytes and record exist.

\subsection{Function Artifact Directory Convention}

Each function should use a predictable internal hierarchy. Finance is the first
implementation and should use:

\begin{verbatim}
business_containers/<business_id>/finance/
    spreadsheets/
    bank-statements/
    invoices-receipts/
    management-accounts/
    vat-tax/
    payroll/
    payment-processors/
    crm-sales/
    reports/
    models/
    evidence-receipts/
\end{verbatim}

Other functions should use equivalent domain-specific folders. A function must
declare its folder taxonomy in one schema/configuration module so the uploader,
artifact service, File Cabinet, and renderer do not each invent names.

The original extension must be preserved. For example, \texttt{.xlsx},
\texttt{.xls}, \texttt{.csv}, \texttt{.pdf}, \texttt{.ofx},
\texttt{.qif}, \texttt{.xml}, and \texttt{.json} remain distinguishable.
File classification determines the folder; the filename alone must not determine
whether the content is trusted or parsed.

\subsection{File Cabinet Contract}

The backend-authoritative File Cabinet tree is stored at:

\begin{verbatim}
.runtime/AION_BUSINESS/workflow_file_cabinets/<business_id>/tree.json
\end{verbatim}

Its repository and API are:

\begin{verbatim}
backend/modules/aion_business/runtime/workflow_file_cabinet_repository.py
backend/modules/aion_business/api/workflow_file_cabinet_api.py
GET  /api/aion/business/workflow-file-cabinet/<business_id>
POST /api/aion/business/workflow-file-cabinet/<business_id>
\end{verbatim}

The File Cabinet must contain pointer nodes, not copies of the source data. The
existing pointer type is \texttt{business\_container\_artifact}, with
\texttt{source\_of\_truth} set to \texttt{business\_container} and
\texttt{file\_cabinet\_role} set to \texttt{index\_pointer\_only}.

Every function must ensure its top-level department folder and required
subfolders exist. Finance should appear as:

\begin{verbatim}
Finance
    Spreadsheets
    Bank Statements
    Invoices & Receipts
    Management Accounts
    VAT & Tax
    Payroll
    Payment Processors
    CRM & Sales
    Reports
    Financial Models
    Evidence & Receipts
\end{verbatim}

A pointer should include:

\begin{verbatim}
{
  "id": "finance_pointer_...",
  "name": "cashflow-2026.xlsx",
  "type": "business_container_artifact",
  "document_type": "spreadsheet",
  "source_of_truth": "business_container",
  "file_cabinet_role": "index_pointer_only",
  "target": {
    "business_container_id": "home-fixed",
    "sub_container": "finance",
    "artifact_type": "spreadsheet",
    "storage_path": "business_containers/home-fixed/finance/spreadsheets/...",
    "artifact_id": "artifact_...",
    "artifact_hash": "sha256:...",
    "receipt_hash": "sha256:...",
    "replay_locator_hash": "sha256:..."
  }
}
\end{verbatim}

Deleting a File Cabinet pointer must not silently delete the canonical artifact.
Artifact deletion or archival requires a separate guarded operation and audit
receipt.

\subsection{Document Ingestion and Evidence Extraction}

Upload, persistence, parsing, and acceptance as business truth are separate
operations:

\begin{enumerate}
  \item \textbf{Upload}: persist original bytes and calculate a hash.
  \item \textbf{Index}: create the File Cabinet pointer.
  \item \textbf{Parse}: extract cells, text, transactions, or metadata.
  \item \textbf{Map}: propose Business Map facts and function-model fields.
  \item \textbf{Validate}: apply schema, period, currency, and consistency checks.
  \item \textbf{Confirm}: obtain user confirmation where confidence or meaning
  is uncertain.
  \item \textbf{Promote}: write accepted facts/model revisions into canonical
  containers.
\end{enumerate}

Selecting a file must not imply permission to write to an external accounting
system. Parsing a file must not imply that every extracted number is correct.

\subsection{Connected Systems}

Accounting, CRM, payment, calendar, email, and other connectors must separate:

\begin{itemize}
  \item discovered or selected;
  \item authentication pending;
  \item read permission granted;
  \item initial import completed;
  \item synchronisation active;
  \item write permission requested;
  \item write permission approved for a bounded action; and
  \item disconnected or revoked.
\end{itemize}

Read access never grants write access. Write access never grants authority for
every future action. Payments, filings, customer communications, bookings,
publishing, and production changes require specific approval policies.

\subsection{Boardroom Context Envelope}

The Boardroom should receive a slim, structured, immutable data sheet generated
from canonical Business Containers. It should not send every raw file or rely on
provider memory.

A representative envelope is:

\begin{verbatim}
{
  "schema_version": "aion.boardroom_context_envelope.v1",
  "business_id": "home-fixed",
  "snapshot_id": "boardroom_context_...",
  "generated_at": "...",
  "business_identity": {},
  "business_structure": {},
  "business_map_projection": {
    "verified_facts": [],
    "material_assumptions": [],
    "unknowns": [],
    "conflicts": []
  },
  "function_models": {
    "finance": {},
    "marketing": {},
    "sales": {},
    "operations": {},
    "support": {}
  },
  "department_intelligence": {},
  "goals": {},
  "active_packages": {},
  "latest_results": {},
  "alerts_and_risks": [],
  "evidence_index": [],
  "previous_boardroom_decisions": [],
  "permissions": {},
  "source_revisions": {},
  "context_hash": "sha256:..."
}
\end{verbatim}

Every participating provider receives the same envelope and meeting instruction.
Provider-specific formatting may adapt syntax but must not alter facts. The
provider call record should retain the envelope hash, provider, model, prompt
version, response hash, timestamps, and safety configuration.

The existing Boardroom session packet already records an execution boundary in
which approval is required, live external side effects are false, raw model tool
access is false, and provider memory mutation is false. These boundaries must be
preserved in the canonical implementation.

\subsection{Boardroom Output and Consensus Persistence}

Boardroom responses are proposals until consensus and user approval are recorded.
The canonical meeting artifact should contain:

\begin{itemize}
  \item the exact context-envelope hash;
  \item council members, providers, and models;
  \item each independent response;
  \item critiques and conflicts;
  \item consensus result;
  \item user amendments;
  \item final approved decisions;
  \item function assignments;
  \item approval boundaries;
  \item evidence requirements;
  \item previous-session references; and
  \item artifact, receipt, and replay hashes.
\end{itemize}

Approved decisions become versioned function packages. They must reference the
Boardroom session and context snapshot that produced them. Functions execute
against the active package and return results, evidence, and exceptions to their
ledger for the next Boardroom projection.

\subsection{Cross-Function Data Flow}

The standard lifecycle for every function is:

\begin{verbatim}
Business Foundation
    -> Canonical Foundation commit
    -> Business Map facts
    -> Function setup and source discovery
    -> Function discovery draft
    -> Function model
    -> Department Intelligence summary
    -> Boardroom context projection
    -> Boardroom decision and approved package
    -> Function execution
    -> Results, evidence and receipts
    -> Business Map/model updates
    -> Next Boardroom context projection
\end{verbatim}

The Business Twin Orchestrator owns progression between functions. A function
runtime must not directly unlock the next department, activate the Boardroom, or
claim a package has been approved.

\subsection{Required Backend APIs}

The platform requires guarded APIs for:

\begin{itemize}
  \item registering and resolving canonical business identity;
  \item committing the Foundation packet;
  \item reading and appending Business Map facts;
  \item reading and updating function discovery drafts;
  \item reading and versioning function models;
  \item reading and updating Department Intelligence records;
  \item uploading and committing artifacts;
  \item parsing artifacts as a separate operation;
  \item reading and writing File Cabinet pointers;
  \item generating a Boardroom Context Envelope;
  \item committing Boardroom sessions and decisions; and
  \item recording approvals, receipts, and replay locators.
\end{itemize}

Every mutation response should include \texttt{ok}, business identifier,
container kind, record or revision identifier, updated timestamp, hash, receipt,
and sync state. Failures must leave the previous canonical record intact.

\subsection{Current Technical Gaps to Resolve}

The code inspection identified these current gaps:

\begin{enumerate}
  \item the active Home Fixed runtime has an operational summary but does not yet
  contain canonical Business Identity or Business Structure files produced from
  onboarding;
  \item Foundation and Business Map data are still principally frontend packets;
  \item function discovery and Department Intelligence are principally stored in
  browser storage;
  \item no canonical Finance model container exists yet;
  \item the artifact helper constructs a committed-looking record but does not
  physically persist artifact bytes;
  \item the File Cabinet is backend-authoritative but the active Home Fixed tree
  can be empty and requires deterministic function-folder provisioning;
  \item frontend and backend business-ID normalisation is inconsistent;
  \item the frontend Boardroom session packet reads browser state while the
  backend Boardroom reads canonical container projections; and
  \item a single canonical Boardroom Context Envelope has not yet replaced these
  parallel context paths.
\end{enumerate}

These gaps must be fixed centrally rather than separately inside each function.

\subsection{Implementation Order}

The recommended implementation order is:

\begin{enumerate}
  \item create and enforce the canonical business-ID resolver;
  \item migrate or alias legacy business IDs without losing data;
  \item add a guarded Foundation-to-Business-Container commit;
  \item promote the Business Map to a canonical backend contract;
  \item register the function-model and function-ledger container contracts;
  \item implement Finance as the first full vertical slice;
  \item implement physical artifact upload and receipt persistence;
  \item implement deterministic File Cabinet folder provisioning and pointer
  insertion;
  \item make browser storage a cache with explicit sync status;
  \item implement the canonical Boardroom Context Envelope;
  \item route every provider the same immutable context snapshot;
  \item persist Boardroom decisions and versioned function packages; and
  \item generalise the completed Finance pattern into shared function modules.
\end{enumerate}

\subsection{Shared Function Module Boundaries}

To avoid further growth of the monolithic desktop application file, reusable
logic should be divided into modules such as:

\begin{verbatim}
desktop/mac/src/business_twin/
    canonical_business_identity.js
    business_container_client.js
    business_map_client.js
    function_ledger_client.js
    artifact_client.js
    file_cabinet_pointer_client.js
    boardroom_context_client.js

desktop/mac/src/departments/shared/
    function_discovery_runtime.js
    function_setup_renderer.js
    function_source_registry.js
    function_model_lifecycle.js

desktop/mac/src/departments/<function>/
    <function>_discovery_schema.js
    <function>_model_schema.js
    <function>_pilot_runtime.js
    <function>_workspace.js
\end{verbatim}

The main application renderer should only select the active function and delegate
to its module. It must not contain each function's schemas, persistence logic,
upload handlers, or model calculations.

\subsection{Finance Pilot Implementation Baseline}

The Finance implementation establishes the reusable baseline for subsequent
functions. The desktop modules are
\texttt{aion\_finance\_discovery\_schema.js},
\texttt{aion\_finance\_model\_builder.js}, and
\texttt{aion\_finance\_pilot.js}. Canonical synchronisation and operational
queries pass through \texttt{aion\_business\_container\_client.js}; lifecycle
progression remains owned by
\texttt{aion\_business\_twin\_orchestrator.js}.

Finance discovery records reporting currency, pricing, monthly revenue, sales
volume, revenue pattern, payment terms, direct costs, fixed overheads, owner
labour and capacity, cash and buffer, receivables, payables, tax, debt, targets,
and approval limits. Answers are founder-provided facts until reconciled to
evidence. The model builder calculates annualised revenue, gross profit, gross
margin, operating surplus, and operating margin only when the required numeric
inputs exist. Every metric stores its formula and is marked
\texttt{calculated\_from\_unverified\_input}; missing inputs and parsing
assumptions remain explicit.

On finalisation, Finance emits a completion receipt to the Business Twin
Orchestrator. The orchestrator marks Finance complete, marks Operations
\texttt{ready\_to\_activate}, and records Operations as the next department.
Finance itself neither routes to nor activates Operations. The permanent
Finance terminal sends read-only questions to
\texttt{/api/aion/business/data/finance-agent/\{business\_id\}/turn}. Answers
are grounded in the canonical Business Financial Model, remain approval-gated,
and are appended to its bounded operational transcript.

\subsection{Safety and Approval Rules}

The following boundaries apply across every function:

\begin{itemize}
  \item discovery does not grant execution authority;
  \item file selection does not grant external-system access;
  \item read access does not grant write access;
  \item a model recommendation does not approve an action;
  \item Boardroom consensus does not replace required user approval;
  \item provider responses cannot mutate canonical truth directly;
  \item public posting, messaging, spending, payments, bookings, filings,
  accounting writes, deployments, and destructive actions require explicit
  bounded authority;
  \item every approved mutation requires evidence and a receipt; and
  \item permissions must be revalidated at execution time.
\end{itemize}

\subsection{Testing and Acceptance Requirements}

Every function implementation must include tests proving:

\begin{enumerate}
  \item all records use the active canonical business identifier;
  \item no global, temporary, or demo container can receive persistent data;
  \item discovery resumes after restart without duplicate canonical writes;
  \item frontend cache and backend truth reconcile correctly;
  \item facts retain source, confidence, evidence, and timestamps;
  \item model revisions preserve previous versions or supersession references;
  \item uploaded files remain inside the active business/function directory;
  \item traversal filenames and unsupported types are rejected;
  \item artifact hashes and receipts reproduce correctly;
  \item File Cabinet rows are pointers to existing canonical artifacts;
  \item File Cabinet deletion does not delete canonical bytes implicitly;
  \item Department Intelligence references the current function-model revision;
  \item Boardroom context contains bounded structured data, not arbitrary raw
  files;
  \item every provider receives the same context hash;
  \item provider memory mutation remains disabled;
  \item unapproved external writes are blocked; and
  \item the new function does not directly unlock another function.
\end{enumerate}

An end-to-end acceptance test should create a new business, complete Foundation,
upload function evidence, finish function discovery, construct a model, verify
the Business Container files, verify File Cabinet pointers, generate a Boardroom
Context Envelope, conduct a provider fan-out using one context hash, approve a
package, execute a safe internal task, and confirm that results and receipts
appear in the next Boardroom projection.

\subsection{Reusable Function Build Checklist}

Before declaring any new function complete, confirm:

\begin{itemize}
  \item[$\square$] The function uses the canonical business-ID resolver.
  \item[$\square$] Its discovery schema is modular and versioned.
  \item[$\square$] Its draft is resumable and backend-synchronised.
  \item[$\square$] Its facts write to the Business Map with provenance.
  \item[$\square$] Its specialist model has a canonical container contract.
  \item[$\square$] Its Department Intelligence entry is persistent.
  \item[$\square$] Its model and ledger have distinct responsibilities.
  \item[$\square$] Its evidence uses canonical artifact persistence.
  \item[$\square$] Its File Cabinet hierarchy and pointer mapping are defined.
  \item[$\square$] Its Boardroom summary is bounded and data-only.
  \item[$\square$] Its context is included in the shared Boardroom envelope.
  \item[$\square$] Its provider calls use the common context hash.
  \item[$\square$] Its plans and actions are approval-gated.
  \item[$\square$] Its results, evidence, and receipts feed back into the
  Business Container.
  \item[$\square$] Its runtime does not own cross-function orchestration.
  \item[$\square$] Its implementation does not add substantive logic to the
  monolithic main application renderer.
  \item[$\square$] Restart, reset, migration, and full new-business tests pass.
\end{itemize}

\subsection{Definition of Done}

The architecture is complete when AION can reconstruct the complete current
business state from the scoped Business Container without relying on browser
storage or provider memory; locate every user-facing artifact through File
Cabinet pointers; produce a compact, hashed Boardroom Context Envelope; deliver
the identical envelope to every participating Boardroom provider; persist the
approved decision and function package; and feed subsequent function results back
into the next business-state projection with full provenance, evidence, receipts,
and approval history.
